\documentclass[10pt,a4paper]{amsart}
\usepackage{amsmath,amssymb,amsthm,graphicx}
\usepackage{verbatim}
\newenvironment{reference}{\comment}{\endcomment}
%\newenvironment{reference}{}{}
\newenvironment{slogan}{\comment}{\endcomment}
\newenvironment{history}{\comment}{\endcomment}

% For commutative diagrams we use Xy-pic
\usepackage[all]{xy}

% We use 2cell for 2-commutative diagrams.
\xyoption{2cell}
% Theorem environments.
%
\theoremstyle{plain}
\newtheorem{theorem}[subsection]{Theorem}
\newtheorem{proposition}[subsection]{Proposition}
\newtheorem{lemma}[subsection]{Lemma}

\theoremstyle{definition}
\newtheorem{definition}[subsection]{Definition}
\newtheorem{example}[subsection]{Example}
\newtheorem{exercise}[subsection]{Exercise}
\newtheorem{situation}[subsection]{Situation}

\theoremstyle{remark}
\newtheorem{remark}[subsection]{Remark}
\newtheorem{remarks}[subsection]{Remarks}

\numberwithin{equation}{subsection}

% Macros
%
\def\lim{\mathop{\mathrm{lim}}\nolimits}
\def\colim{\mathop{\mathrm{colim}}\nolimits}
\def\Spec{\mathop{\mathrm{Spec}}}
\def\Hom{\mathop{\mathrm{Hom}}\nolimits}
\def\Ext{\mathop{\mathrm{Ext}}\nolimits}
\def\SheafHom{\mathop{\mathcal{H}\!\mathit{om}}\nolimits}
\def\SheafExt{\mathop{\mathcal{E}\!\mathit{xt}}\nolimits}
\def\Sch{\mathit{Sch}}
\def\Mor{\mathop{\mathrm{Mor}}\nolimits}
\def\Ob{\mathop{\mathrm{Ob}}\nolimits}
\def\Sh{\mathop{\mathit{Sh}}\nolimits}
\def\NL{\mathop{N\!L}\nolimits}
\def\CH{\mathop{\mathrm{CH}}\nolimits}
\def\proetale{{pro\text{-}\acute{e}tale}}
\def\etale{{\acute{e}tale}}
\def\QCoh{\mathit{QCoh}}
\def\Ker{\mathop{\mathrm{Ker}}}
\def\Im{\mathop{\mathrm{Im}}}
\def\Coker{\mathop{\mathrm{Coker}}}
\def\Coim{\mathop{\mathrm{Coim}}}

% Boxtimes
%
\DeclareMathSymbol{\boxtimes}{\mathbin}{AMSa}{"02}

%
% Macros for moduli stacks/spaces
%
\def\QCohstack{\mathcal{QC}\!\mathit{oh}}
\def\Cohstack{\mathcal{C}\!\mathit{oh}}
\def\Spacesstack{\mathcal{S}\!\mathit{paces}}
\def\Quotfunctor{\mathrm{Quot}}
\def\Hilbfunctor{\mathrm{Hilb}}
\def\Curvesstack{\mathcal{C}\!\mathit{urves}}
\def\Polarizedstack{\mathcal{P}\!\mathit{olarized}}
\def\Complexesstack{\mathcal{C}\!\mathit{omplexes}}
% \Pic is the operator that assigns to X its picard group, usage \Pic(X)
% \Picardstack_{X/B} denotes the Picard stack of X over B
% \Picardfunctor_{X/B} denotes the Picard functor of X over B
\def\Pic{\mathop{\mathrm{Pic}}\nolimits}
\def\Picardstack{\mathcal{P}\!\mathit{ic}}
\def\Picardfunctor{\mathrm{Pic}}
\def\Deformationcategory{\mathcal{D}\!\mathit{ef}}
\usepackage{cleveref}
\crefname{originallemma}{original lemma}{original lemmas}
\Crefname{originallemma}{Original lemma}{Original lemmas}
\makeatletter
\newlabel{descent-lemma-local-source-target-characterize}{{04R1}{source}{}{sourcepointer-1}{}}
\newlabel{descent-lemma-local-source-target-characterize@cref}{{[originallemma][0][]04R1}{[1][0][]source}{}{}{}}
\makeatother
\crefname{originalsection}{original section}{original sections}
\Crefname{originalsection}{Original section}{Original sections}
\makeatletter
\newlabel{section-representable}{{03HA}{source}{}{sourcepointer-2}{}}
\newlabel{section-representable@cref}{{[originalsection][0][]03HA}{[1][0][]source}{}{}{}}
\makeatother
\crefname{originallemma}{original lemma}{original lemmas}
\Crefname{originallemma}{Original lemma}{Original lemmas}
\makeatletter
\newlabel{spaces-lemma-lift-morphism-presentations}{{02X1}{source}{}{sourcepointer-3}{}}
\newlabel{spaces-lemma-lift-morphism-presentations@cref}{{[originallemma][0][]02X1}{[1][0][]source}{}{}{}}
\makeatother
\title{Testing an etale local on the source-and-target property of morphisms on algebraic spaces}
\author{The Stacks Project Authors}
\date{stacks-project:spaces-morphisms.tex:local-snapshot-20261012; GNU FDL 1.2 or later}
\begin{document}
\maketitle
\setcounter{subsection}{1}
\noindent\textit{Editorial note.} All delivered mathematical text is the author's own, copied from the pinned Stacks Project chapter file spaces-morphisms.tex; no statement or proof step was rewritten, reordered or removed.\par\medskip
\noindent\textit{Editorial note.} The statement is a meta-statement about an arbitrary property P; the proof handles each of the three equivalence claims in its own paragraph, and the closing paragraph on representable f is part of the same proof environment.\par\medskip
\noindent\textit{Editorial note.} The three references of the counted spans are recorded as non-load-bearing navigation pointers; the argument that consumes them stays inside the counted proof span.\par\medskip
\noindent\textit{Editorial note.} Preamble: the three native excerpt spans of sources/stacks-spaces-morphisms/preamble.tex (lines 8-13, 14-18 and 154-217). The chapter file itself carries no preamble text.\par\medskip
\noindent\textit{Allowed standard prerequisites (editorial).}\begin{itemize}
\item Etale localness on the source-and-target characterized through a double diagram. Descent, Lemma descent-lemma-local-source-target-characterize, used at proof lines 4372-4376 and 4379-4382. Is the engine of both the (2) implies (1) direction and the representable comparison.
\item Lifting a morphism from an algebraic space to a presentation by a scheme. Spaces, Lemma spaces-lemma-lift-morphism-presentations, used at proof lines 4343-4345. Makes (1) implies (2) immediate by exhibiting a diagram with U -> X surjective.
\item The definition of a property of representable morphisms in the functorial sense. Section section-representable, used at proof lines 4336-4339. Fixes the third equivalent formulation in the statement.
\item Base change and fppf locality of the property. used at proof lines 4385-4390. Allows the final comparison for representable f.
\end{itemize}\medskip
\ifdefined\hypertarget\hypertarget{sourcepointer-1}{}\fi
\noindent\textit{Source reference.} Original lemma 04R1: focused-seeds/AGP/sources/stacks-snapshot-20261011/tags.txt, source lines 6169--6169. Navigation pointer into the Stacks Project; the standard fact it supplies is applied inside the counted proof span and the argument resting on it stays there.\par\medskip
\ifdefined\hypertarget\hypertarget{sourcepointer-2}{}\fi
\noindent\textit{Source reference.} Original section 03HA: focused-seeds/AGP/sources/stacks-snapshot-20261011/tags.txt, source lines 4559--4559. Navigation pointer into the Stacks Project; the standard fact it supplies is applied inside the counted proof span and the argument resting on it stays there.\par\medskip
\ifdefined\hypertarget\hypertarget{sourcepointer-3}{}\fi
\noindent\textit{Source reference.} Original lemma 02X1: focused-seeds/AGP/sources/stacks-snapshot-20261011/tags.txt, source lines 3824--3824. Navigation pointer into the Stacks Project; the standard fact it supplies is applied inside the counted proof span and the argument resting on it stays there.\par\medskip
\begin{lemma}
\label{lemma-local-source-target}
Let $\mathcal{P}$ be a property of morphisms of schemes
which is \'etale local on the source-and-target.
Let $S$ be a scheme.
Let $f : X \to Y$ be a morphism of algebraic spaces over $S$.
Consider commutative diagrams
$$
\xymatrix{
U \ar[d]_a \ar[r]_h & V \ar[d]^b \\
X \ar[r]^f & Y
}
$$
where $U$ and $V$ are schemes and the vertical arrows are \'etale.
The following are equivalent
\begin{enumerate}
\item for any diagram as above the morphism $h$ has property
$\mathcal{P}$, and
\item for some diagram as above with $a : U \to X$ surjective
the morphism $h$ has property $\mathcal{P}$.
\end{enumerate}
If $X$ and $Y$ are representable, then this is also equivalent to
$f$ (as a morphism of schemes) having property $\mathcal{P}$.
If $\mathcal{P}$ is also preserved under any base change, and
fppf local on the base, then for representable morphisms $f$ this
is also equivalent to $f$ having property $\mathcal{P}$ in the sense
of Section \ref{section-representable}.
\end{lemma}

\begin{proof}
Let us prove the equivalence of (1) and (2).
The implication (1) $\Rightarrow$ (2) is immediate (taking into account
Spaces, Lemma \ref{spaces-lemma-lift-morphism-presentations}).
Assume
$$
\xymatrix{
U \ar[d] \ar[r]_h & V \ar[d] \\
X \ar[r]^f & Y
}
\quad\quad
\xymatrix{
U' \ar[d] \ar[r]_{h'} & V' \ar[d] \\
X \ar[r]^f & Y
}
$$
are two diagrams as in the lemma. Assume $U \to X$ is
surjective and $h$ has property $\mathcal{P}$. To show that (2)
implies (1) we have to prove that $h'$ has $\mathcal{P}$. To do
this consider the diagram
$$
\xymatrix{
U \ar[d]_h &
U \times_X U' \ar[l] \ar[d]^{(h, h')} \ar[r] &
U' \ar[d]^{h'} \\
V &
V \times_Y V' \ar[l] \ar[r] &
V'
}
$$
By
Descent, Lemma \ref{descent-lemma-local-source-target-characterize}
we see that $h$ has $\mathcal{P}$ implies $(h, h')$ has $\mathcal{P}$
and since $U \times_X U' \to U'$ is surjective this implies (by the
same lemma) that $h'$ has $\mathcal{P}$.

\medskip\noindent
If $X$ and $Y$ are representable, then
Descent, Lemma \ref{descent-lemma-local-source-target-characterize}
applies which shows that (1) and (2) are equivalent to $f$ having
$\mathcal{P}$.

\medskip\noindent
Finally, suppose $f$ is representable, and $U, V, a, b, h$ are
as in part (2) of the lemma, and that $\mathcal{P}$ is preserved under
arbitrary base change. We have to show that for any scheme
$Z$ and morphism $Z \to X$ the base change $Z \times_Y X \to Z$
has property $\mathcal{P}$. Consider the diagram
$$
\xymatrix{
Z \times_Y U \ar[d] \ar[r] &
Z \times_Y V \ar[d] \\
Z \times_Y X \ar[r] &
Z
}
$$
Note that the top horizontal arrow is a base change of $h$ and
hence has property $\mathcal{P}$. The left vertical arrow is \'etale
and surjective and the right vertical arrow is \'etale. Thus
Descent, Lemma \ref{descent-lemma-local-source-target-characterize}
once again kicks in and shows that $Z \times_Y X \to Z$
has property $\mathcal{P}$.
\end{proof}

\end{document}
